%------------------------------------------------------------------------------%
\begin{question}
  Calcule
  \(
    \int \frac{x-3}{x(x-1)(x+1)} \,dx
  \)
\end{question}

%------------------------------------------------------------------------------%
\begin{resolution}{Solução}
  O grau do numerador
  \[
    P(x) = x - 2
  \]
  \pause
  é menor do que o do denominador
  \[
    Q(x) = x(x-1)(x+1)
  \]
  \pause
  O denominador já está fatorado

  \pause
  Por frações parciais temos
  \[
    \pause
    f(x)
    \;=\; \frac{P(x)}{Q(x)}
    \pause
    \;=\; \frac{x-3}{x(x-1)(x+1)}
    \pause
    \;=\; \frac{A}{x} + \frac{B}{x-1} + \frac{C}{x+1}
  \]
\end{resolution}

%------------------------------------------------------------------------------%
\begin{resolution}{Solução}
  \begin{align*}
    \onslide<+->{\frac{x-3}{x(x-1)(x+1)}}
    \onslide<+->{& = \frac{A}{x}
                   + \frac{B}{x-1}
                   + \frac{C}{x+1}
    \\[1mm]}
    \onslide<+->{x-3
                 & = \left(
                     \frac{A}{x}
                   + \frac{B}{x-1}
                   + \frac{C}{x+1}
                   \right)
                   x(x-1)(x+1)
    \\[1mm]}
    \onslide<+->{& = A(x-1)(x+1)
                   + Bx(x+1)
                   + Cx(x-1)
    \\}
    \onslide<+->{& = A\left(x^2 + 1\right)
                   + B\left(x^2 + x\right)
                   + C\left(x^2 - x\right)
    \\}
    \onslide<+->{& = x^2 \left( A + B + C \right)
                   + x   \left(B-C\right)
                   - A}
  \end{align*}
\end{resolution}

%------------------------------------------------------------------------------%
\begin{resolution}{Solução}
  \[
    \onslide<+->{
    x - 3
    \;=\; x^2 \left(A + B + C \right)
    \;+\; x   \left(B - C\right)
    \;-\; A}
  \]

\vspace{-\baselineskip}
\begin{columns}[t]
\column{0.3\linewidth}
  \[
    \onslide<+->{\systeme[sort={A,B,C}]{
      A + B + C = 0,
          B - C = 1,
      A         = 3
    }}
  \]

  \smallskip
  \[
    \onslide<+->{B = 1 + C}
  \]
\column{0.4\linewidth}
  \begin{align*}
    \onslide<+->{A + B + C     & = 0  \\}
    \onslide<+->{3 + 1 + C + C & = 0  \\}
    \onslide<+->{2C            & = -4 \\}
    \onslide<+->{C             & = -2}
  \end{align*}
  \[
    \onslide<+->{B = 1 - 2 = -1}
  \]
\end{columns}
\end{resolution}

%------------------------------------------------------------------------------%
\begin{resolution}{Solução}
\begin{align*}
  \onslide<+->{\frac{x-3}{x(x-1)(x+1)}
  & = \frac{A}{x}
    + \frac{B}{x-1}
    + \frac{C}{x+1} \\[2mm]}
  \onslide<+->{& = \frac{3}{x}
    + \frac{-1}{x-1}
    + \frac{-2}{x+1} \\[2mm]}
  \onslide<+->{& = \frac{3}{x}
    - \frac{1}{x-1}
    - \frac{2}{x+1}}
\end{align*}
\end{resolution}

%------------------------------------------------------------------------------%
\begin{resolution}{Solução}
\begin{align*}
  \onslide<+->{\int \frac{x-3}{x(x-1)(x+1)} \,dx}
  \onslide<+->{& = \int \frac{3}{x}
    - \frac{1}{x-1}
    - \frac{2}{x+1} \,dx \\[2mm]}
  \onslide<+->{& = \int \frac{3}{x} \,dx
    - \int \frac{1}{x-1} \,dx
    - 2 \int \frac{1}{x+1} \,dx \\[2mm]}
  \onslide<+->{& = 3\ln\abs{x}
    - \ln\abs{x-1}
    - 2 \ln\abs{x+1}
    + k \\[2mm]}
  \onslide<+->{&\color{gray}
    = \ln\abs{\frac{x^3}{(x-1)(x+1)^2}}
    + k}
\end{align*}
\end{resolution}

%------------------------------------------------------------------------------%
